% ============================================================================
% Fichier  : partie1/P01_C02_controle_interne.tex
% Titre    : Le Contrôle Interne et ses Composantes
% Auteur   : MINKA MI NGUIDJOÏ Thierry Emmanuel
% Version  : 1.0 -- octobre 2026
% ============================================================================
% Sources : COSO, Internal Control -- Integrated Framework (2013), cinq
% composantes et dix-sept principes paraphrasés ; règlement COBAC R-2016/04,
% articles 1 et 2 (texte relevé sur une base réglementaire en ligne pour le
% cours GCI, à collationner avec le texte officiel) ; articles A0 et A1 pour
% la lecture µ, U, H.
% ============================================================================

\chapter{Le Contrôle Interne et ses Composantes}
\label{chap:controle-interne}

\epigraph{%
    « Un contrôle qui ne laisse pas de trace\\
    n'existe que dans la procédure qui le décrit. »%
}{Principe de l'auditeur}

\niveauChapitre{Fondamental}{%
    Maîtrise du Chapitre~\ref{chap:audit-auditologie} (audit, contrôle,
    trois lignes)%
}

\begin{boxobjectifs}
\begin{itemize}
    \item Définir le contrôle interne et ses trois catégories d'objectifs.
    \item Énoncer les cinq composantes et les dix-sept principes du
          COSO~2013.
    \item Classer un contrôle~: préventif ou détectif, manuel ou automatisé,
          général ou d'application.
    \item Appliquer la séparation des tâches et la règle des quatre yeux à
          une matrice des droits.
    \item Concevoir un dispositif de contrôle interne à l'aide d'une matrice
          des risques et des contrôles.
    \item Évaluer la conception et l'efficacité d'un contrôle, et coter le
          résultat \OUI{}, \NON{} ou \NEK{}.
\end{itemize}
\end{boxobjectifs}

\bigskip

% ============================================================================
\section*{Question directrice : où passaient les 186 millions~?}
% ============================================================================

L'écart de stock de LiZbiZ (Chapitre~\ref{chap:prologue}) ne s'explique
pas par un seul acte. Les premières investigations d'AUDIX CONSEIL montrent
que, dans l'application intégrée, un même magasinier peut enregistrer une
réception de marchandises \emph{et} ajuster le stock après inventaire. Le
responsable informatique dispose de tous les droits, y compris celui de
valider un paiement. Personne ne revoit ces droits.

Aucune règle n'interdisait ces cumuls~: LiZbiZ n'a jamais conçu son
contrôle interne, il s'est formé au fil des urgences. Ce chapitre définit
le contrôle interne (section~\ref{sec:c02-definition}), en présente le cadre
de référence (section~\ref{sec:c02-coso}), classe les contrôles
(section~\ref{sec:c02-types}), applique la séparation des tâches
(section~\ref{sec:c02-sod}), puis montre comment concevoir et évaluer un
dispositif (sections~\ref{sec:c02-concevoir} et~\ref{sec:c02-evaluer}).

\separateur

% ============================================================================
\section{Qu'est-ce que le Contrôle Interne~?}
\label{sec:c02-definition}
% ============================================================================

\subsection{Une définition}

Le cadre de référence le plus répandu est celui du \sigle{COSO}
(\textit{Committee of Sponsoring Organizations of the Treadway Commission}),
dont la version actuelle date de 2013~\cite{coso2013ic}.

\begin{boxdef}
\termecle{Contrôle interne}\index{contrôle interne}~: processus mis en œuvre
par l'organe de gouvernance, la direction et l'ensemble du personnel, destiné
à fournir une assurance raisonnable quant à la réalisation d'objectifs liés
aux opérations, à l'information (\textit{reporting}) et à la conformité.
\end{boxdef}
{\small\itshape Définition reformulée d'après le COSO~2013~\cite{coso2013ic}.}

\medskip

Quatre idées sont contenues dans cette définition.
\begin{enumerate}
    \item \textbf{Un processus}, pas un document~: le contrôle interne vit
          dans les opérations quotidiennes.
    \item \textbf{L'affaire de tous}~: il est porté par la direction et le
          personnel, pas seulement par le contrôleur interne.
    \item \textbf{Une assurance raisonnable}, jamais absolue
          (Chapitre~\ref{chap:audit-auditologie}).
    \item \textbf{Des objectifs}~: on ne contrôle pas « en général », mais en
          vue d'objectifs précis.
\end{enumerate}

\subsection{Trois catégories d'objectifs}

\begin{tableaucours}{L{2.8cm}L{5.0cm}L{4.4cm}}{%
    \tableauHeader{Catégorie} &
    \tableauHeader{Question} &
    \tableauHeader{Exemple chez LiZbiZ}
}
Opérations & Les activités sont-elles efficaces et efficientes, les actifs
protégés~? & Le stock enregistré correspond au stock réel \\
Information & L'information financière et non financière est-elle
fiable~? & Les états financiers reflètent la valeur du stock \\
Conformité & Les lois et règlements sont-ils respectés~? & Les données
personnelles des clients sont protégées (loi n°~2024/017) \\
\end{tableaucours}
\vspace{-0.8em}
\begin{center}{\small\itshape Tableau~2.1. Les objectifs du contrôle
interne}\end{center}

% ============================================================================
\section{Le Cadre COSO 2013}
\label{sec:c02-coso}
% ============================================================================

\subsection{Cinq composantes, dix-sept principes}

Le COSO organise le contrôle interne en \textbf{cinq composantes}, détaillées
en \textbf{dix-sept principes}. Un dispositif n'est efficace que si les cinq
composantes sont présentes, fonctionnent et agissent ensemble.

\begin{tableaucours}{L{3.2cm}cL{8.0cm}}{%
    \tableauHeader{Composante} &
    \tableauHeader{N°} &
    \tableauHeader{Principe (reformulé)}
}
Environnement de contrôle & 1 & Afficher un engagement envers l'intégrité et
l'éthique \\
 & 2 & Faire superviser le contrôle interne par un organe de gouvernance
indépendant de la direction \\
 & 3 & Définir les structures, les rattachements, les pouvoirs et les
responsabilités \\
 & 4 & Attirer, développer et retenir des personnes compétentes \\
 & 5 & Faire rendre compte chacun de ses responsabilités de contrôle \\
Évaluation des risques & 6 & Fixer des objectifs suffisamment clairs pour
identifier les risques \\
 & 7 & Identifier et analyser les risques qui menacent ces objectifs \\
 & 8 & Prendre en compte le risque de fraude \\
 & 9 & Identifier et analyser les changements significatifs \\
Activités de contrôle & 10 & Choisir et mettre en place des activités de
contrôle qui réduisent les risques \\
 & 11 & Choisir et mettre en place des contrôles généraux sur les
technologies \\
 & 12 & Déployer les contrôles par des politiques et des procédures \\
Information et communication & 13 & Produire et utiliser une information
pertinente et de qualité \\
 & 14 & Communiquer en interne sur le contrôle interne \\
 & 15 & Communiquer avec les tiers \\
Pilotage & 16 & Conduire des évaluations continues ou ponctuelles \\
 & 17 & Évaluer les déficiences et les communiquer à temps \\
\end{tableaucours}
\vspace{-0.8em}
\begin{center}{\small\itshape Tableau~2.2. Les composantes et les principes
du COSO~2013 (notre reformulation)}\end{center}

\begin{boxinfo}
\textbf{Le principe~11 concerne directement l'auditeur des SI.} Le COSO
demande de choisir et de mettre en place des contrôles généraux sur les
technologies~: accès, changements, exploitation. C'est le cœur des
contrôles généraux informatiques (section~\ref{sec:c02-types}).
\end{boxinfo}

\subsection{Le diagnostic de LiZbiZ}

\begin{boxscenario}{Opération LIZBIZ. Les cinq composantes}
\textbf{Constatations} (entretiens et documents remis à AUDIX CONSEIL)~:
\begin{itemize}
    \item \emph{Environnement de contrôle}~: pas de comité d'audit~; les
          pouvoirs de validation ne sont écrits nulle part.
    \item \emph{Évaluation des risques}~: aucune cartographie des risques~;
          le risque de fraude n'a jamais été examiné.
    \item \emph{Activités de contrôle}~: des contrôles existent dans les
          agences, mais varient d'un site à l'autre.
    \item \emph{Information et communication}~: les écarts d'inventaire ne
          remontent à la direction qu'une fois par an.
    \item \emph{Pilotage}~: aucune évaluation du dispositif depuis la
          création de l'entreprise.
\end{itemize}

\textbf{Analyse} (confiance moyenne). Les activités de contrôle ne peuvent
pas compenser l'absence des quatre autres composantes. L'écart de stock est
le symptôme d'un dispositif qui n'a jamais été conçu.
\end{boxscenario}

% ============================================================================
\section{Classer les Contrôles}
\label{sec:c02-types}
% ============================================================================

\subsection{Selon le moment}

\begin{boxdef}
\termecle{Contrôle préventif}\index{contrôle préventif}~: empêche
l'anomalie de se produire (une validation obligatoire avant paiement).

\medskip
\termecle{Contrôle détectif}\index{contrôle détectif}~: repère l'anomalie
après coup (le rapprochement mensuel du stock et des inventaires).

\medskip
\termecle{Contrôle correctif}\index{contrôle correctif}~: corrige l'anomalie
et ses effets (la restauration d'une sauvegarde après une suppression).
\end{boxdef}

\subsection{Selon le mode d'exécution}

Un contrôle peut être \textbf{manuel} (une signature sur un bon de sortie),
\textbf{automatisé} (l'application refuse un ajustement de stock non
validé) ou \textbf{semi-automatisé} (l'application produit une liste
d'anomalies qu'une personne examine). Un contrôle automatisé est fiable et
constant, à une condition~: que personne ne puisse le désactiver ou le
modifier sans contrôle. D'où la distinction suivante.

\subsection{Contrôles généraux et contrôles d'application}

\begin{boxdef}
\termecle{Contrôles généraux informatiques}\index{contrôles généraux
informatiques} (\textit{IT general controls})~: contrôles qui s'appliquent à
l'environnement dans lequel les applications sont développées, exploitées et
maintenues~: gestion des accès, gestion des changements, exploitation
(sauvegardes, traitements, incidents).

\medskip
\termecle{Contrôles d'application}\index{contrôles d'application}~:
contrôles intégrés à une application, qui agissent sur le traitement des
opérations~: contrôles des entrées (validation des données), des traitements
(exactitude des calculs) et des sorties (complétude des états).
\end{boxdef}

\begin{boxalert}
\textbf{Violation typique.} L'auditeur constate que l'application refuse tout
paiement d'un montant supérieur au plafond autorisé, et conclut que le
contrôle est efficace. Or le responsable informatique peut modifier le
plafond directement dans la base, sans trace.

\smallskip
\textbf{Règle~:} un contrôle d'application ne vaut que si les contrôles
généraux qui le protègent (accès, changements) sont eux-mêmes efficaces.
L'auditeur les teste en premier.
\end{boxalert}

% ============================================================================
\section{Séparer les Tâches}
\label{sec:c02-sod}
% ============================================================================

\subsection{Le principe}

\begin{boxdef}
\termecle{Séparation des tâches}\index{séparation des tâches}
(\textit{segregation of duties}, \sigle{SoD})~: répartition, entre des
personnes différentes, des fonctions dont le cumul permettrait de commettre
et de dissimuler une erreur ou une fraude. On sépare en principe quatre
fonctions~: \emph{autoriser} une opération, l'\emph{enregistrer},
\emph{garder} l'actif concerné, et \emph{contrôler}.

\medskip
\termecle{Règle des quatre yeux}\index{règle des quatre yeux}~: toute
opération sensible est réalisée par une personne et validée par une autre.
\end{boxdef}

\subsection{La matrice des droits de LiZbiZ}

\begin{boxscenario}{Opération LIZBIZ. Qui peut faire quoi~?}
AUDIX CONSEIL extrait de l'application intégrée les droits de six
utilisateurs sur cinq fonctions sensibles~: F1 créer ou modifier un
fournisseur~; F2 saisir une commande d'achat~; F3 enregistrer une réception
en entrepôt~; F4 ajuster le stock après inventaire~; F5 valider un paiement
fournisseur.

\medskip
\begin{tableaucours}{L{4.6cm}ccccc}{%
    \tableauHeader{Utilisateur (fonction)} &
    \tableauHeader{F1} & \tableauHeader{F2} & \tableauHeader{F3} &
    \tableauHeader{F4} & \tableauHeader{F5}
}
M.~KOUNOU (resp. informatique) & × & × & × & × & × \\
M.~MBIDA (acheteur) & × & × & & & \\
M.~FOUDA (magasinier, Douala) & & & × & × & \\
Mme~OWONA (comptable) & × & & & & × \\
M.~NANGA (chef d'agence, Garoua) & & × & × & & \\
Mme~MBALLA (DAF) & & & & & × \\
\end{tableaucours}

\medskip
\textbf{Paires incompatibles retenues.} F1 et F5 (créer un fournisseur
fictif, puis le payer)~; F2 et F3 (commander, puis attester soi-même une
réception qui n'a pas eu lieu)~; F3 et F4 (enregistrer une réception, puis
« corriger » le stock pour masquer un manquant)~; F2 et F5 (commander, puis
payer sa propre commande).

\medskip
\textbf{Constatation.} Quatre utilisateurs sur six cumulent au moins une
paire incompatible~: M.~KOUNOU (toutes), M.~FOUDA (F3 et F4), Mme~OWONA
(F1 et F5) et M.~NANGA (F2 et F3).

\medskip
\textbf{Analyse} (confiance élevée). Le cumul F3--F4 du magasinier de Douala
suffit à expliquer qu'un manquant puisse disparaître des écritures. Le
cumul total du responsable informatique est le plus grave~: il cumule les
droits de gestion et le pouvoir technique de modifier les droits eux-mêmes.
\end{boxscenario}

\subsection{Quand on ne peut pas séparer}

Dans une petite agence de trois personnes, une séparation complète est
impossible. On recourt alors à des \termecle{contrôles
compensatoires}\index{contrôle compensatoire}~: revue a posteriori par une
personne extérieure à l'agence, rapprochement mensuel par le siège, journal
des ajustements examiné par le contrôleur interne. Le contrôle compensatoire
doit être écrit, exécuté et tracé, comme tout autre contrôle.

% ============================================================================
\section{Concevoir un Dispositif}
\label{sec:c02-concevoir}
% ============================================================================

Mettre en place un contrôle interne, c'est relier chaque risque à un ou
plusieurs contrôles, et organiser la preuve de leur exécution.

\begin{enumerate}
    \item \textbf{Fixer les objectifs} du processus (par exemple~: le stock
          enregistré est exact et complet).
    \item \textbf{Identifier les risques} qui menacent ces objectifs
          (Chapitre~\ref{chap:risque}).
    \item \textbf{Décrire le processus} et ses points de risque.
    \item \textbf{Choisir les contrôles}~: préventifs d'abord, détectifs en
          complément, automatisés lorsque c'est possible.
    \item \textbf{Documenter} chaque contrôle dans une matrice des risques
          et des contrôles.
    \item \textbf{Tester} puis \textbf{piloter} le dispositif.
\end{enumerate}

\begin{boxprocedure}
\textbf{La matrice des risques et des contrôles.} Pour que le dispositif soit
auditable, chaque ligne de la matrice indique~:
\begin{itemize}
    \item le processus, l'objectif et le risque traité~;
    \item le contrôle, sa nature (préventif ou détectif, manuel ou
          automatisé) et sa fréquence~;
    \item le responsable de l'exécution~;
    \item \textbf{la preuve} que laisse chaque exécution (signature,
          journal, état signé)~;
    \item la date de la dernière évaluation et son résultat.
\end{itemize}
Un contrôle sans preuve d'exécution ne peut pas être évalué~: il doit être
reconçu.
\end{boxprocedure}

\begin{boxmetier}{Contrôleur interne}
\textbf{Mettre en place le contrôle des stocks chez LiZbiZ}
\begin{enumerate}
    \item Retirer à M.~FOUDA le droit F4~; confier les ajustements de stock
          à une personne qui ne gère pas les réceptions.
    \item Rendre l'ajustement soumis à validation (quatre yeux) au-delà d'un
          seuil, par exemple 500~000~FCFA.
    \item Retirer au responsable informatique tous les droits de gestion~;
          tracer ses accès d'administration.
    \item Instaurer un inventaire tournant mensuel dans chaque entrepôt, et
          un rapprochement par le siège.
    \item Consigner le tout dans la matrice des risques et des contrôles, et
          prévoir une revue trimestrielle des droits.
\end{enumerate}
\end{boxmetier}

% ============================================================================
\section{Évaluer un Contrôle}
\label{sec:c02-evaluer}
% ============================================================================

\subsection{Conception et efficacité}

L'auditeur évalue un contrôle en deux temps.

\begin{tableaucours}{L{3.2cm}L{4.6cm}L{4.4cm}}{%
    \tableauHeader{Évaluation} &
    \tableauHeader{Question} &
    \tableauHeader{Comment}
}
Conception & Le contrôle, s'il fonctionne, réduit-il le risque~? &
Entretien, lecture de la procédure, test de cheminement sur un cas \\
Efficacité opérationnelle & Le contrôle a-t-il fonctionné comme prévu
pendant toute la période~? & Test d'un échantillon d'exécutions et de leurs
preuves \\
\end{tableaucours}
\vspace{-0.8em}
\begin{center}{\small\itshape Tableau~2.3. Deux niveaux d'évaluation d'un
contrôle}\end{center}

Un contrôle mal conçu est inefficace, même s'il est exécuté à la lettre. Un
contrôle bien conçu mais non exécuté ne protège de rien.

\subsection{\NON{} ou \NEK{}~?}

La cotation d'un contrôle demande une distinction que l'audit binaire ne
fait pas.

\begin{itemize}
    \item \textbf{\NON{}}~: l'auditeur établit que le contrôle n'est pas
          conforme au critère. Par exemple, le critère exige une validation
          tracée de chaque ajustement, et l'extraction montre des
          ajustements sans validation.
    \item \textbf{\NEK{}}~: l'information disponible ne permet pas de
          conclure, malgré des diligences normales. Par exemple, le journal
          des ajustements n'a pas été remis, ou deux sources se contredisent
          sans qu'on puisse les départager.
\end{itemize}

\begin{boxalert}
\textbf{Violation typique.} Coter \NEK{} un contrôle dont le critère exige une
trace, alors que l'auditeur a vérifié qu'aucune trace n'existe.

\smallskip
\textbf{Règle~:} l'absence d'une trace exigée par le critère est un \NON{}
établi. \NEK{} est réservé au cas où l'auditeur ne peut pas savoir~; il
s'accompagne toujours de la description des diligences effectuées.
\end{boxalert}

\subsection{Lire un domaine}

\begin{boxlecture}{Contrôles des stocks de LiZbiZ}
Douze contrôles de la gestion des stocks sont évalués sur la période de
janvier à avril~2026~: 1~\OUI{}, 10~\NON{}, 1~\NEK{}.

\medskip
\textbf{Constatation.} $\mu = 1/12 \approx 0{,}08$~; $U = 1/12 \approx 0{,}08$~;
$H \approx 0{,}82$~bit.

\medskip
\textbf{Analyse} (confiance élevée). La conformité est très faible, mais
l'incertitude l'est aussi~: l'organisation sait que ses contrôles de stock
ne fonctionnent pas. C'est une \emph{non-conformité connue}, qui appelle des
correctifs directs, à l'inverse du domaine~D du prologue ($H \approx 1{,}24$),
où il fallait d'abord comprendre.

\medskip
\hfill\lecture{0{,}08}{0{,}08}{0{,}82}
\end{boxlecture}

\subsection{Les limites du contrôle interne}

Même bien conçu, un dispositif ne donne qu'une assurance raisonnable. Il
peut être mis en échec par la collusion entre plusieurs personnes, par un
dirigeant qui passe outre, par une erreur de jugement, ou parce que le coût
d'un contrôle dépasserait le risque qu'il traite. L'auditeur en tient compte
en recherchant, au-delà des contrôles, les indices de leur contournement
(Chapitre~\ref{chap:fraude}).

% ============================================================================
\section{Le Contrôle Interne dans les Banques}
\label{sec:c02-cobac}
% ============================================================================

Pour les établissements de crédit de la CEMAC, le contrôle interne n'est
pas une option~: c'est une obligation réglementaire~\cite{cobac2016r04}.

\begin{boxjuris}
\textbf{Règlement COBAC R-2016/04, articles~1 et~2} (texte relevé sur une
base réglementaire en ligne, à collationner avec le texte officiel)~:

\medskip
Article~1~: \textit{« Le présent règlement fixe les exigences minimales
relatives au système de contrôle interne que les établissements assujettis
doivent mettre en œuvre. »}

\medskip
Article~2, définition du risque opérationnel~: \textit{« risque de pertes
résultant de l'inadéquation ou de la défaillance des processus internes,
des personnes et des systèmes, ou d'évènements externes. »}

\medskip
\textbf{Analyse.} La défaillance d'un système informatique est, par
définition, un risque opérationnel que le dispositif de contrôle interne de
la banque doit maîtriser. L'auditeur des SI d'une banque teste donc les
contrôles informatiques au titre de ce règlement, et pas seulement au titre
des bonnes pratiques.
\end{boxjuris}

\begin{boxscenario}{Opération SANAGA. Le contrôle interne après la panne}
\textbf{Faits.} Après la panne de trois jours de mars~2026, la COBAC demande
à la BANQUE SANAGA de démontrer que son dispositif de contrôle interne
couvre les risques liés aux systèmes d'information.

\medskip
\textbf{Constatation} (audit interne, Mme~EKAMBI). La matrice des risques
et des contrôles de la banque compte 214~contrôles, dont 9~seulement portent
sur l'informatique, et aucun sur la gestion des changements.

\medskip
\textbf{Analyse} (confiance élevée). Le risque qui s'est réalisé, un
changement non maîtrisé du système central, n'était couvert par aucun
contrôle~: le défaut est de conception, pas d'exécution. Le
Chapitre~\ref{chap:audit-continuite} reviendra sur l'audit de la continuité
de la banque.
\end{boxscenario}

\separateur

\begin{boxresume}
\textbf{Ce qu'il faut retenir de ce chapitre~:}
\begin{itemize}
    \item \textbf{Le contrôle interne est un processus}, porté par tous,
          qui donne une assurance raisonnable sur les opérations,
          l'information et la conformité.
    \item \textbf{Cinq composantes, dix-sept principes} (COSO~2013)~: les
          activités de contrôle ne suffisent pas sans environnement,
          évaluation des risques, information et pilotage.
    \item \textbf{Un contrôle se classe} selon le moment (préventif,
          détectif, correctif), le mode (manuel, automatisé) et la portée
          (général ou d'application)~; les contrôles généraux conditionnent
          les contrôles d'application.
    \item \textbf{La séparation des tâches} interdit le cumul des fonctions
          qui permettraient de commettre et de dissimuler~; à défaut, des
          contrôles compensatoires tracés.
    \item \textbf{Un dispositif se conçoit} dans une matrice des risques et
          des contrôles, où chaque contrôle laisse une preuve.
    \item \textbf{Évaluer un contrôle}, c'est juger sa conception puis son
          efficacité~; l'absence d'une trace exigée est un \NON{}, pas un
          \NEK{}.
\end{itemize}
\end{boxresume}

\bigskip

\begin{boxexo}[title={\ding{45}\quad Exercice~2.1. Classer des contrôles \niveaubase}]
Classez chacun des contrôles suivants (préventif, détectif ou correctif~;
manuel ou automatisé~; général ou d'application).
\begin{enumerate}[(a)]
    \item L'application refuse une facture dont le numéro de commande
          n'existe pas.
    \item Le responsable informatique revoit chaque trimestre la liste des
          comptes à privilèges.
    \item La comptable rapproche chaque mois le relevé bancaire et le grand
          livre.
    \item Une sauvegarde de la base est restaurée chaque mois sur un serveur
          de test.
\end{enumerate}
\end{boxexo}

\begin{boxexo}[title={\ding{45}\quad Exercice~2.2. Corriger la matrice des droits \niveaumoyen}]
Reprenez la matrice des droits de LiZbiZ (section~\ref{sec:c02-sod}).
\begin{enumerate}[(a)]
    \item Proposez une nouvelle répartition des droits qui supprime toutes
          les paires incompatibles, sans retirer à quiconque les droits
          nécessaires à son métier.
    \item L'agence de Garoua ne compte que trois personnes. Proposez deux
          contrôles compensatoires pour le cumul F2--F3 de M.~NANGA.
    \item Rédigez la ligne de la matrice des risques et des contrôles
          correspondant à l'un de ces contrôles.
\end{enumerate}
\end{boxexo}

\begin{boxexo}[title={\ding{45}\quad Exercice~2.3. Le dispositif de la BANQUE SANAGA \niveauexpert}]
La COBAC demande à la BANQUE SANAGA de compléter son dispositif de contrôle
interne sur la gestion des changements informatiques.
\begin{enumerate}[(a)]
    \item Identifiez trois risques liés aux changements du système central.
    \item Pour chacun, concevez un contrôle, en précisant sa nature, sa
          fréquence, son responsable et la preuve qu'il laisse.
    \item Pour l'un de ces contrôles, rédigez le test de conception et le
          test d'efficacité que réalisera l'audit interne.
    \item Indiquez dans quel cas l'auditrice devra coter le contrôle \NEK{}
          plutôt que \NON{}.
\end{enumerate}
\end{boxexo}

\bigskip

\begin{boxinfo}
\textbf{Transition.} Un contrôle se justifie par le risque qu'il traite. Le
Chapitre~\ref{chap:risque} définit le risque, apprend à le mesurer et à le
représenter dans une cartographie, point de départ de tout dispositif de
contrôle et de tout plan d'audit.
\end{boxinfo}
